\documentclass[a4paper,12pt]{article}
\usepackage[utf8]{inputenc}
\usepackage[T1]{fontenc}
\usepackage[ngerman]{babel}
\usepackage{amsmath, amssymb}
\usepackage{geometry}
\usepackage{tikz}
\usetikzlibrary{shapes.geometric, arrows.meta, positioning, fit, backgrounds}
\usepackage{parskip}
\usepackage{titlesec}
\usepackage{booktabs}
\usepackage{enumitem}

\geometry{left=3cm, right=3cm, top=2.5cm, bottom=2.5cm}

\titleformat{\section}{\large\bfseries}{}{0em}{}[\titlerule]
\titleformat{\subsection}{\normalsize\bfseries}{}{0em}{}

\tikzset{
  box/.style={rectangle, draw, rounded corners=3pt, text width=4.2cm,
              align=center, minimum height=0.9cm, font=\small, inner sep=4pt},
  boxw/.style={box, text width=5.2cm},
  boxn/.style={box, text width=3.8cm},
  arrow/.style={-{Latex[length=2mm]}, thick},
}

\begin{document}

\begin{center}
  {\LARGE\bfseries Studierhinweise zu Lektion 4}\\[0.5em]
  {\large Analysis (Modul 61211)}
\end{center}

\vspace{1em}

Diese Kurseinheit ist vornehmlich drei zentralen Sätzen der Analysis gewidmet:
dem lokalen Umkehrsatz, dem Satz über implizite Funktionen und dem Satz von
Taylor für Funktionen mehrerer Veränderlicher. Die beiden ersten Sätze gelten
als ,,schwierig'' -- konzentrieren Sie sich auf das Verständnis ihrer Aussagen
und deren Bedeutung für die Auflösung nichtlinearer Gleichungssysteme. Die
Beweise können Sie zunächst übergehen; wichtig ist das darin verwendete
Kontraktionsprinzip (Banachscher Fixpunktsatz), das in weiten Teilen der
Mathematik eine Rolle spielt.

Im dritten Abschnitt werden Kenntnisse aus MG wiederholt: Ableitungen höherer
Ordnung, Polynom- und Potenzreihenfunktionen, Taylorpolynom und Satz von
Taylor. Neu ist vielleicht die Leibnizsche Formel für $(fg)^{(k)}$, die eine
genaue Analogie zur binomischen Formel darstellt. Der vierte Abschnitt führt
partielle und totale Differenzierbarkeit höherer Ordnung für Funktionen
mehrerer Veränderlicher ein und erweitert den Satz von Taylor auf diesen Fall.

Der letzte Abschnitt ist der Extremwertbestimmung gewidmet. Die erste
Ableitung liefert wie im eindimensionalen Fall eine notwendige Bedingung;
hinreichende Bedingungen erhält man nun über die Definitheit der Hesseschen
Matrix der partiellen Ableitungen zweiter Ordnung. Lokale Extrema unter
Nebenbedingungen werden mit der Methode der Lagrangemultiplikatoren bestimmt.

% -----------------------------------------------------------------------
\section*{Struktur der Kurseinheit 4}

\begin{center}
\begin{tikzpicture}[node distance=0.5cm and 0.6cm]

  % Column headers
  \node[font=\bfseries\small] (hA) at (0,0)    {Rückblick auf MG};
  \node[font=\bfseries\small] (hB) at (9.5,0)  {Ergänzungen und neue Inhalte};

  % ---- Left column: MG-Rückblick ----
  \node[box, below=0.8cm of hA] (k1) {4.1 $\mathrm{Hom}(\mathbb{R}^n, \mathbb{R}^n)$\\und Matrizen};
  \node[box, below=0.5cm of k1] (k2) {4.2 $Az = 0$,\\lin.\ Gl.-System mit\\$(n, m+n)$-Matrix $A$};
  \node[box, below=0.5cm of k2] (k3) {4.3 $k$-te Ableitung für\\$f: M \to \mathbb{R}$ ($M \subseteq \mathbb{R}$)};
  \node[box, below=0.5cm of k3] (k4) {4.3 Taylorpolynom,\\Satz von Taylor};
  \node[box, below=0.5cm of k4] (k5) {4.5 lok.\ Extremum,\\notw.\ Bedingung};
  \node[box, below=0.5cm of k5] (k6) {4.5 hinreichende\\Bedingung};

  % ---- Right column: new content ----
  \node[box, below=0.8cm of hB] (n1) {4.1 Umkehrsatz für\\$f \in \mathrm{Hom}(\mathbb{R}^n, \mathbb{R}^n)$};
  \node[box, below=0.5cm of n1] (n2) {4.1 Banachscher\\Fixpunktsatz};
  \node[box, below=0.5cm of n2] (n3) {4.1 Invertierbarkeit\\benachbarter Matrizen};
  \node[box, below=0.5cm of n3] (n4) {4.1 lokaler Umkehrsatz};
  \node[box, below=0.5cm of n4] (n5) {4.2 Satz über\\implizite Funktionen};
  \node[box, below=0.5cm of n5] (n6) {4.4 $k$-malige partielle,\\$k$-malige (stetige)\\Differenzierbarkeit};
  \node[box, below=0.5cm of n6] (n7) {4.4 Schwarzscher\\Vertauschbarkeitssatz};
  \node[box, below=0.5cm of n7] (n8) {4.4 Taylorpolynom,\\Satz von Taylor};
  \node[box, below=0.5cm of n8] (n9) {4.5 lokales Extremum,\\notwendige Bedingung};
  \node[box, below=0.5cm of n9] (n10){4.5 hinreichende Bedingung:\\Hessesche Matrix};
  \node[box, below=0.5cm of n10] (n11){4.5 Extremum unter\\Nebenbedingungen,\\Lagrangemultiplikatoren};

  % Arrows left column
  \draw[arrow] (k1) -- (k2);
  \draw[arrow] (k2) -- (k3);
  \draw[arrow] (k3) -- (k4);
  \draw[arrow] (k4) -- (k5);
  \draw[arrow] (k5) -- (k6);

  % Arrows right column
  \draw[arrow] (n1) -- (n2);
  \draw[arrow] (n2) -- (n3);
  \draw[arrow] (n3) -- (n4);
  \draw[arrow] (n4) -- (n5);
  \draw[arrow] (n5) -- (n6);
  \draw[arrow] (n6) -- (n7);
  \draw[arrow] (n7) -- (n8);
  \draw[arrow] (n8) -- (n9);
  \draw[arrow] (n9) -- (n10);
  \draw[arrow] (n10) -- (n11);

  % Cross arrows (left -> right)
  \draw[arrow, dashed] (k1.east) -- ++(0.3,0) |- (n1.west);
  \draw[arrow, dashed] (k2.east) -- ++(0.3,0) |- (n5.west);
  \draw[arrow, dashed] (k4.east) -- ++(0.3,0) |- (n8.west);
  \draw[arrow, dashed] (k6.east) -- ++(0.3,0) |- (n10.west);

\end{tikzpicture}
\end{center}

\newpage

% -----------------------------------------------------------------------
\section*{Zielelement 4.1 -- Der Umkehrsatz}

\subsection*{Lerninhalte}

\begin{center}
\begin{tikzpicture}[node distance=0.5cm and 1.0cm]
  \node[boxw] (umk)  {Umkehrsatz für $f \in \mathrm{Hom}(\mathbb{R}^n, \mathbb{R}^n)$:\\eindeutige Lösbarkeit von $Ax = y$\\mit quadratischer Matrix $A$};
  \node[boxw, below=0.5cm of umk] (lok) {lokaler Umkehrsatz:\\lokale eindeutige Lösbarkeit von\\$f(x) = y$ mit $f \in \mathrm{Abb}(M, \mathbb{R}^n)$, $M \subseteq \mathbb{R}^n$};
  \node[boxw, below=0.5cm of lok] (inv) {Lösbarkeit benachbarter\\linearer Gleichungssysteme\\($A_{nn}$, Stetigkeit von $\mathrm{inv}$)};
  \node[boxw, below=0.5cm of inv] (ban) {Banachscher Fixpunktsatz\\(Kontraktionsprinzip)};
  \node[boxw, below=0.5cm of ban] (off) {offene Abbildung};
  \draw[arrow] (umk) -- (lok);
  \draw[arrow] (lok) -- (inv);
  \draw[arrow] (inv) -- (ban);
  \draw[arrow] (ban) -- (off);
\end{tikzpicture}
\end{center}

\subsection*{Lernziele}

Nach Durcharbeiten dieses Abschnitts sollten Sie sich aus MG wieder vergegenwärtigt haben,
\begin{itemize}
  \item welche Bedeutung der Rang einer Matrix für die Lösbarkeit eines linearen
        Gleichungssystems hat,
  \item wie man mithilfe der Treppennormalform den Rang (und ggf.\ die Inverse)
        einer Matrix bestimmen kann,
\end{itemize}
und Sie sollten in der Lage sein,
\begin{itemize}
  \item den lokalen Umkehrsatz zu formulieren und auf konkrete Beispiele anzuwenden,
  \item den Banachschen Fixpunktsatz zu formulieren und anzuwenden.
\end{itemize}

\subsection*{Selbstkontrollelement 4.1}

Bleibt der Banachsche Fixpunktsatz richtig, wenn auf die Abgeschlossenheit der
Menge $M$ verzichtet wird?

\newpage

% -----------------------------------------------------------------------
\section*{Zielelement 4.2 -- Implizit definierte Funktionen}

\subsection*{Lerninhalte}

\begin{center}
\begin{tikzpicture}[node distance=0.5cm and 1.0cm]
  \node[boxw] (lin)  {Satz über implizite lineare Funktionen:\\Lösbarkeit und Lösung von\\$Az = 0$ mit $(n, m+n)$-Matrix $A$};
  \node[boxw, below=0.5cm of lin] (imp) {Satz über implizite Funktionen:\\lokale Auflösbarkeit von\\$F(z) = 0$ mit $F \in \mathrm{Abb}(M, \mathbb{R}^n)$, $M \subseteq \mathbb{R}^{m+n}$};
  \node[boxw, below=0.5cm of imp] (abl) {Ableitung der impliziten Funktion:\\$g'(x) = -[F_y(x, g(x))]^{-1} F_x(x, g(x))$};
  \draw[arrow] (lin) -- (imp);
  \draw[arrow] (imp) -- (abl);
\end{tikzpicture}
\end{center}

\subsection*{Lernziele}

Nach Durcharbeiten dieses Abschnitts sollten Sie in der Lage sein,
\begin{itemize}
  \item den Satz über implizite Funktionen zu formulieren und auf konkrete Beispiele anzuwenden.
\end{itemize}

\subsection*{Selbstkontrollelement 4.2}

Sei $F: \mathbb{R}^2 \to \mathbb{R}$, $(x, y) \mapsto F(x, y) := x^3 - xy + y^3$ gegeben.
Die Abbildung stellt die Menge $B := \{(x, y) \in \mathbb{R}^2 \mid F(x, y) = 0\}$
(,,Kartesisches Blatt'') dar. Können Sie
\begin{itemize}
  \item zum Punkt $(\tfrac{1}{2}, 1)$ eine Umgebung einzeichnen, in der $F(x, y) = 0$ sowohl
        eindeutig nach $y$ als auch eindeutig nach $x$ aufgelöst werden kann,
  \item alle Punkte einzeichnen, in deren Umgebung es keine Auflösung von $F(x, y) = 0$
        (i)~nach $y$, (ii)~nach $x$ gibt?
\end{itemize}

\newpage

% -----------------------------------------------------------------------
\section*{Zielelement 4.3 -- Rückblick und Ergänzungen: Ableitungen höherer Ordnung}

\subsection*{Lerninhalte}

\begin{center}
\begin{tikzpicture}[node distance=0.5cm and 1.0cm]
  \node[boxw] (abl)  {Ableitung $k$-ter Ordnung\\für $f: M \to \mathbb{R}$, $M \subseteq \mathbb{R}$:\\$f^{(0)} := f$, $f^{(k+1)} := (f^{(k)})'$};
  \node[boxw, below=0.5cm of abl] (op) {Operationen: Summe, Produkt,\\Quotient, Verkettung;\\\textbf{Leibnizsche Formel}};
  \node[boxw, below=0.5cm of op] (pot) {Anwendungen:\\Polynomfunktionen,\\Potenzreihenfunktionen,\\Identitätssatz};
  \node[boxw, below=0.5cm of pot] (kp) {$k$-te Ableitung in einem Punkt $a$};
  \node[boxw, below=0.5cm of kp] (tp) {Taylorpolynom $T_k$ mit\\Entwicklungspunkt $a$};
  \node[boxw, below=0.5cm of tp] (tay) {Satz von Taylor:\\Restfunktion $R_k$,\\Restdarstellung von Lagrange};
  \draw[arrow] (abl) -- (op);
  \draw[arrow] (op) -- (pot);
  \draw[arrow] (pot) -- (kp);
  \draw[arrow] (kp) -- (tp);
  \draw[arrow] (tp) -- (tay);
\end{tikzpicture}
\end{center}

\subsection*{Lernziele}

Nach Durcharbeiten dieses Abschnitts sollten Sie sich Gelerntes aus MG wieder
vergegenwärtigt haben, im Einzelnen sollten Sie (wieder) in der Lage sein,
\begin{itemize}
  \item zu erklären, wann eine Funktion $f$ in einem Punkt $a$ ihres Definitionsbereichs
        $k$-mal differenzierbar ist,
  \item das Taylorpolynom $T_k$ von $f$ mit Entwicklungspunkt $a$ hinzuschreiben,
  \item den Satz von Taylor (mit der Restdarstellung von Lagrange) zu formulieren
        und im Hinblick auf die Approximierbarkeit von $f$ durch $T_k$ zu interpretieren.
\end{itemize}

\subsection*{Selbstkontrollelement 4.3}

Sie sollten nachweisen können (gegebenenfalls nach Konsultation der Lösung von
Ü~3.3.4), dass $f: \mathbb{R} \to \mathbb{R}$ mit
\[
  f(x) := \begin{cases} x^3 \sin\bigl(\tfrac{1}{x}\bigr), & \text{falls } x \neq 0, \\ 0, & \text{falls } x = 0, \end{cases}
\]
differenzierbar ist und eine stetige Ableitung besitzt, aber nicht zweimal
differenzierbar ist.

\newpage

% -----------------------------------------------------------------------
\section*{Zielelement 4.4 -- Ableitungen höherer Ordnung (mehrere Veränderliche)}

\subsection*{Lerninhalte}

\begin{center}
\begin{tikzpicture}[node distance=0.5cm and 1.0cm]
  \node[boxw] (par)  {$k$-malige partielle\\Differenzierbarkeit,\\partielle Ableitungen\\der Ordnung $k$};
  \node[boxw, below=0.5cm of par] (sch) {Schwarzscher\\Vertauschbarkeitssatz};
  \node[boxw, below=0.5cm of sch] (dif) {$k$-malige (totale) Differenzierbarkeit,\\$k$-malige stetige Differenzierbarkeit\\$f \in C^k(M)$};
  \node[boxw, below=0.5cm of dif] (tp) {Taylorpolynom $T_k$ für\\$f: M \to \mathbb{R}$, $M \subseteq \mathbb{R}^n$};
  \node[boxw, below=0.5cm of tp] (tay) {Satz von Taylor,\\Approximationsgeschwindigkeit};
  \draw[arrow] (par) -- (sch);
  \draw[arrow] (sch) -- (dif);
  \draw[arrow] (dif) -- (tp);
  \draw[arrow] (tp) -- (tay);
\end{tikzpicture}
\end{center}

\subsection*{Lernziele}

Nach Durcharbeiten dieses Abschnitts sollten Sie in der Lage sein,
\begin{itemize}
  \item die verschiedenen Begriffe der Differenzierbarkeit höherer Ordnung zu formulieren
        und zueinander in Beziehung zu setzen,
  \item den Satz von Taylor zu formulieren und auf konkret gegebene Funktionen anzuwenden,
        insbesondere die Taylorpolynome einer Funktion zu gegebenem Entwicklungspunkt
        aufzustellen.
\end{itemize}

\subsection*{Selbstkontrollelement 4.4}

Ist Ihnen klar geworden, dass der Mittelwertsatz~3.5.9 (unter der schärferen
Voraussetzung, dass $f$ auf $[a, b]$ differenzierbar ist und $[a, b] \subseteq M$
gilt) als Spezialfall im Satz von Taylor~4.4.8 enthalten ist?

\newpage

% -----------------------------------------------------------------------
\section*{Zielelement 4.5 -- Extrema}

\subsection*{Lerninhalte}

\begin{center}
\begin{tikzpicture}[node distance=0.5cm and 1.0cm]
  \node[boxw] (ext)  {lokales Extremum in einem Punkt $a$\\(Maximum, Minimum)};
  \node[boxw, below=0.5cm of ext] (not) {notwendige Bedingung:\\$f'(a) = 0$};
  \node[boxw, below=0.5cm of not] (qf) {quadratische Form,\\symmetrische Matrix,\\(semi-)definit, indefinit};
  \node[boxw, below=0.5cm of qf] (hin) {hinreichende Bedingung:\\Hessesche Matrix $H(a)$ in $a$ definit};
  \node[boxw, below=0.5cm of hin] (n2) {explizite Bedingungen im Fall $n = 2$:\\$\Delta(a) = D_{11}D_{22} - D_{12}^2$,\\Sattelpunkt};
  \node[boxw, below=0.5cm of n2] (neb) {lokales Extremum unter\\Nebenbedingungen:\\Lagrangemultiplikatoren};
  \draw[arrow] (ext) -- (not);
  \draw[arrow] (not) -- (qf);
  \draw[arrow] (qf) -- (hin);
  \draw[arrow] (hin) -- (n2);
  \draw[arrow] (n2) -- (neb);
\end{tikzpicture}
\end{center}

\subsection*{Lernziele}

Nach Durcharbeiten dieses Abschnitts sollten Sie in der Lage sein,
\begin{itemize}
  \item den Begriff des lokalen Extremums zu definieren und mit dem des globalen
        Extremums in Zusammenhang zu bringen,
  \item Bedingungen dafür aufzustellen, dass eine zweimal stetig differenzierbare
        Funktion $f: M \to \mathbb{R}$ mit $M \subseteq \mathbb{R}^n$ in einem Punkt ein
        lokales Extremum besitzt,
  \item die Methode der Lagrangemultiplikatoren zur Bestimmung von lokalen
        Extrema unter Nebenbedingungen zu beschreiben.
\end{itemize}

\subsection*{Selbstkontrollelement 4.5}

Können Sie die lokalen Extrema der Funktion $f: \mathbb{R} \to \mathbb{R}$, definiert durch
\[
  f(x) := \begin{cases} |x|, & \text{falls } x \in \mathbb{Q}, \\ 0, & \text{falls } x \in \mathbb{R} \setminus \mathbb{Q}, \end{cases}
\]
bestimmen?

\end{document}
